\documentclass[conference]{IEEEtran}
\IEEEoverridecommandlockouts

\usepackage{amsmath,amssymb}
\usepackage{booktabs}
\usepackage{url}
\usepackage{cite}

\newcommand{\sys}{CoFlow-5 Synapse}

\begin{document}

\title{Single-Writer Multi-Agent Signal Control with a Read-Only Language Layer:\\
A SUMO Case Study of Six Junctions in Tunis Centre-Ville}

% TODO: add the other five team members and the supervisor.
\author{\IEEEauthorblockN{Aziz Messaoud and the CoFlow-5 team}
\IEEEauthorblockA{ESPRIT, Tunis, Tunisia\\
aziz.messaoud@esprit.tn}}

\maketitle

\begin{abstract}
We present \sys{}, a hierarchical multi-agent system that controls six signalised
junctions of Tunis centre-ville in the SUMO microscopic simulator. The design rests on
three rules: each traffic light has exactly one writer, a master agent that arbitrates
proposals from five junction agents in a fixed priority order (safety, emergency,
network directive, late bus, pedestrian, normal); control degrades from Max-Pressure to
actuated to fixed-time when sensor data goes stale; and a language layer explains decisions
from a structured decision log but has no code path to the simulator. Every explanation
must cite log entries, and an answer whose citations do not exist in the log is turned
into an abstention. On a network built from OpenStreetMap with synthetic demand, we compare
the system with SUMO's fixed-time and actuated plans over three scenarios and three random
seeds. The safety rules cut the number of vehicle greens started while pedestrians were
still crossing from 331--405 per run (actuated) to 17--19, and mean pedestrian waiting fell
in two of three scenarios. Vehicle waiting time was mixed ($-24\%$ to $+20\%$ against
actuated), emergency vehicles were slower than with SUMO's plans, and no difference was
statistically significant with three seeds. We report these results as preliminary and use
recent guidance on stochastic benchmarking in SUMO to identify the weaknesses of our own
evaluation protocol, chief among them single-seed parameter tuning and uncorrected multiple
comparisons.
\end{abstract}

\begin{IEEEkeywords}
traffic signal control, multi-agent systems, SUMO, Max-Pressure, large language models,
explainability, reproducibility
\end{IEEEkeywords}

%=====================================================================
\section{Introduction}
%=====================================================================

Most signalised junctions in Tunis centre-ville run fixed-time or locally actuated plans.
Each controller sees only its own queues: it cannot see the road beyond the junction, an
ambulance two junctions away, or a bus running late. Adaptive and learning-based
controllers address part of this, but deploying them raises two practical questions that
are as important as average delay. First, \emph{who is allowed to change a light?} When
several agents (an emergency agent, a transit agent, a learned policy) can each act on the
same signal, their commands can conflict. Second, \emph{why did the light change?} Operators
and incident reviewers need an answer that can be checked, not a fluent guess.

Large language models (LLMs) are increasingly coupled to traffic simulators to build
scenarios and tune parameters from natural language~\cite{li2024chatsumo,jeong2025agentsumo,li2026decoupled}.
In those systems the LLM acts on the simulation. We take the opposite position for signal
control: the language layer may read everything the controllers recorded, but must never
reach the simulator or the signal heads.

This paper makes four contributions:
\begin{enumerate}
\item A hierarchical multi-agent architecture in which a single master per junction is the
only writer of its light, junction agents only propose, and a supermaster coordinates the
network through expiring directives (Section~\ref{sec:design}).
\item A read-only language layer that answers questions from a decision log, cites the log
entries it relies on, and abstains when citations fail a mechanical check
(Section~\ref{sec:synapse}).
\item A reproducible testbed of six real junctions in Tunis built from OpenStreetMap, with
documented limits on lane evidence and demand (Section~\ref{sec:setup}).
\item A preliminary evaluation against SUMO's fixed-time and actuated plans, reported with
its negative results, and a critique of our own protocol against recent benchmarking
guidance~\cite{riehl2026sumoitscontrol} (Sections~\ref{sec:results}--\ref{sec:discussion}).
\end{enumerate}

All results come from simulation with synthetic demand. SUMO emission outputs are proxies,
not air-quality measurements, and simulated emergency trip times are not evidence of lives
saved.

%=====================================================================
\section{Related Work}
%=====================================================================

\subsection{Multi-agent signal control in SUMO}
Agent-based approaches to traffic and transportation are reviewed by Bazzan and
Kl\"ugl~\cite{bazzan2014review}. Azevedo et al.~\cite{azevedo2016jade} coupled the JADE
agent platform to SUMO through the TraSMAPI middleware, binding one agent to each SUMO
traffic light while leaving vehicles as ordinary simulated entities. They contrast two
organisations: one agent per light under a \emph{super-coordinator}, which synchronises
well but has a single point of failure and concentrated computation, and \emph{distributed}
coordination, which synchronises poorly but contains failures. They chose the latter and
had each Q-learning agent combine its own reward with the distance-weighted rewards of its
neighbours. Our design is a hybrid of the two organisations (Section~\ref{sec:design}).
Two other findings of theirs inform our evaluation: a fixed plan with different timings by
time of day beat plain Q-learning on queue waiting time, and waiting time and travel time
ranked the controllers differently.

\subsection{Max-Pressure and classical adaptive control}
Max-Pressure originates in back-pressure scheduling for packet networks~\cite{tassiulas1990stability}
and was introduced to signal control by Varaiya~\cite{varaiya2013maxpressure}, with a
stability guarantee under bounded demand. Kouvelas et al.~\cite{kouvelas2014maxpressure}
added cyclic phase order, fixed cycles and normalised pressure to limit spill-back.
SCOOT~\cite{hunt1981scoot} and SCATS~\cite{lowrie1990scats} adapt cycle length, splits and
offsets for groups of junctions. Riehl et al.~\cite{riehl2026sumoitscontrol} provide
open TraCI implementations of fixed-order and flexible Max-Pressure and a
SCOOT/SCATS-inspired group controller, and argue that SUMO studies too often rely on
single runs and project-specific baselines. They recommend 10--20 replications per
configuration, paired comparisons on identical seeds, variance-aware reporting and formal
hypothesis tests, and show that parameter settings whose means differ by less than one
standard deviation cannot be ranked from a single run~\cite{riehl2026sumoitscontrol,riehl2025reproducibility}.

\subsection{Learning-based control}
Reinforcement learning (RL) for signal control is commonly studied with multi-agent
platforms over SUMO~\cite{bokade2024pytsc}. Our optional RL agent uses Fitted Q-Iteration
with tree ensembles~\cite{ernst2005fqi} and, following~\cite{azevedo2016jade}, a reward that
includes neighbouring junctions.

\subsection{Language models and traffic simulation}
TrafficGPT~\cite{zhang2024trafficgpt}, ChatSUMO~\cite{li2024chatsumo} and
AgentSUMO~\cite{jeong2025agentsumo} connect LLMs to traffic tools so that non-experts can
build and run scenarios. Li and Ke~\cite{li2026decoupled} split the SUMO workflow into
Planner, Builder, Demand, Runner and Analyst agents with an orchestrator that keeps
generated artifacts in a persistent store, so that later agents receive validated file
references instead of relying on conversational memory. Their analyst feeds KPI feedback
back to the planner in the style of Reflexion~\cite{shinn2023reflexion}; one repair round
raised task success on harder tasks from 30\% to 50\%, a second added nothing, and on
easy tasks extra reflection sometimes hurt. Their agents also modify the network and
signal programmes. Our language layer shares the idea of a single validated source of
truth but is deliberately barred from acting.

%=====================================================================
\section{System Design}\label{sec:design}
%=====================================================================

\subsection{Hierarchy}
Table~\ref{tab:hierarchy} lists the agents. Each of the six junctions has one master and
five junction agents. A supermaster runs eight network agents every 5\,s (incident
detection, transit monitoring and emergency routing every 1\,s) and sends directives to the
masters. Agents communicate over a publish/subscribe bus with six message kinds (request,
proposal, directive, status, publication, notification); every message carries an expiry
time, so a silent agent cannot leave a stale request in force.

\begin{table}[t]
\caption{Agents of the hierarchy. Only masters write lights, through the executor.}
\label{tab:hierarchy}
\centering
\footnotesize
\begin{tabular}{@{}p{1.6cm}p{6.4cm}@{}}
\toprule
Level & Agents and role \\
\midrule
Junction ($\times6$) & \textbf{Master}: arbitrates, sole writer of its light. \textbf{Traffic light}: proposes the next phase (Max-Pressure by default; predictive or RL optional). \textbf{Pedestrian}, \textbf{transit priority}, \textbf{emergency priority}: send requests. \textbf{Sensor health}: flags stale data. \\
Network & \textbf{Supermaster}: runs network agents, sends directives, never writes a light. Agents: demand forecast, incident detection, rerouting, emergency route, corridor, transit monitor, weather, notification. \\
Language & \textbf{Synapse}: explainer, operator copilot, alert writer, verification bridge. Read-only. \\
\bottomrule
\end{tabular}
\end{table}

\subsection{Single writer and arbitration}
Junction agents never act on SUMO. They send proposals and requests to their master. Each
simulated second the master resolves pending messages in a fixed order,
\begin{multline}
\text{safety} \succ \text{emergency} \succ \text{directive} \\
\succ \text{late bus} \succ \text{pedestrian} \succ \text{normal},
\end{multline}
and passes at most one phase change to an executor that enforces the safety rules before
writing the light: a minimum green of 5\,s, a 3\,s yellow (longer in rain), a pedestrian
clearance hold of up to 30\,s computed at a design walking speed of 1.07\,m/s, and refusal
of a green whose exit lane is more than 85\% full. Emergency and transit requests are
refused, with a logged reason, when the exit is full, when an earlier emergency vehicle is
being served, or when a bus is not late. A static check run in continuous integration
verifies that only the executor writes lights, that only masters call it, and that
junction agents cannot read scenario events or calendar assumptions.

\subsection{Max-Pressure proposal}
Let $\mathcal{L}(\phi)$ be the vehicle links given right of way by phase $\phi$, with
incoming lane $i(l)$ and outgoing lane $o(l)$; let $\mathcal{C}(\phi)$ be the pedestrian
crossings it serves and $\mathcal{E}(\phi)$ the incoming edges it releases. The traffic
light agent computes
\begin{equation}
p(\phi) = \sum_{l\in\mathcal{L}(\phi)} \big(h_{i(l)} - h_{o(l)}\big)
        + \sum_{e\in\mathcal{E}(\phi)} b_e
        + w_{\text{ped}} \sum_{c\in\mathcal{C}(\phi)} n_c ,
\label{eq:pressure}
\end{equation}
where $h$ is the number of halting vehicles on a lane, $b_e$ a corridor bias sent by the
supermaster for a platoon heading to edge $e$, $n_c$ the number of pedestrians waiting at
crossing $c$, and $w_{\text{ped}}=0.5$. The downstream count $h_{o(l)}$ is taken from the
neighbouring master's status when it is at most 2\,s old, otherwise from local sensors.
Once the current green $\phi_0$ has lasted the minimum green, the agent proposes
$\phi^\star=\arg\max_\phi p(\phi)$ if $p(\phi^\star) > p(\phi_0)+m$, with switching margin
$m=0$. This is the acyclic variant of Max-Pressure~\cite{varaiya2013maxpressure}. Unlike the
flexible controller of~\cite{riehl2026sumoitscontrol} it subtracts downstream queues, but it
currently has no maximum green time (Section~\ref{sec:discussion}).

\subsection{Recovery ladder}
The master selects its mode from data freshness (stale after 10\,s). With all data fresh it
uses Max-Pressure; if downstream data is stale it falls back to actuated control (gap-out
after 3\,s, max-out at 45\,s), which needs incoming lanes only; if incoming data is stale it
runs a fixed cycle. Each change of mode is logged with a reason code. The ladder is our
answer to the single point of failure that~\cite{azevedo2016jade} attributes to
centralised coordination: if the supermaster or a neighbour goes silent, its directives
expire and each junction keeps operating on local data.

\subsection{Learned components behind a rule gate}
Learned models supply estimates, not commands: pedestrian over-waiting, ambulance and bus
arrival times, bus lateness, demand, travel times, rain slowdown and incident detection.
All are scikit-learn~\cite{pedregosa2011sklearn} gradient-boosted trees trained on data
recorded from simulations dated in 2025 and tested on scenarios dated in 2026, so that a
model cannot memorise the date of Ramadan. Each model has a hand-written rule as fallback,
and a model is used only if it beats its rule on held-out data. The optional RL traffic
light agent uses Fitted Q-Iteration~\cite{ernst2005fqi} with $\gamma=0.9$, a reward of own
waiting plus $0.5\times$ the neighbours' mean queue, and a 60\,s starvation guard.

\subsection{Decision log}
Every master and supermaster decision is appended to a log entry
$(\textit{id}, t, \textit{agent}, \textit{event}, \textit{reason}, \textit{tls}, \textit{details})$.
The reason must be one of a fixed table of codes, each with an official plain-language
meaning (for example \texttt{EMERGENCY\_DENIED\_EXIT\_FULL}: ``the exit lane was full, so
green would only move the queue''). Logging an unknown code raises an error, so every
decision has a reason.

%=====================================================================
\section{The Read-Only Language Layer}\label{sec:synapse}
%=====================================================================

Synapse answers questions such as ``Why did J1 hold green at 08:14?'' from a finished or
live decision log. The static architecture check forbids its modules from importing the
simulation backend, sensors, executor or masters, and the operator copilot's
``approve'' action is only recorded, never applied.

\paragraph{Retrieval} Log entries are selected by junction, vehicle and keyword-to-reason-code
matching (for example ``ambulance'' $\rightarrow$ \texttt{EMERGENCY*}, \texttt{RECOVER*}).
No embeddings are used.

\paragraph{Generation} With an API key, a hosted LLM receives only the selected entries and
the meanings of their reason codes. It may call read-only tools (log search, live state,
experiment results) for at most five rounds. It must cite each factual claim as
\texttt{[\#id]} or reply \texttt{ABSTAIN}. Without a key, or if the API fails, a
deterministic template builds the answer from the same entries. The control loop never
waits for the LLM.

\paragraph{Verification} Every cited id is checked against the log. An answer citing an
entry that does not exist, or citing nothing, is rejected and returned as an abstention.
A separate deterministic bridge answers \emph{verify} and \emph{audit} requests: a claim
about a decision is accepted only if it matches the recorded time, action, reason code and
agent, and the answer is labelled \emph{verified}, \emph{template} or \emph{abstain} with
seven provenance fields. Alerts are rewritten into French and Arabic by an LLM, and a
rewrite that changes any number is rejected.

The decision log plays the role that the artifact store plays in~\cite{li2026decoupled}: a
validated record that the language model reads instead of its own memory. The difference is
the direction of authority. In~\cite{li2026decoupled} the LLM's outputs are executed after
validation; here no LLM output is ever executed.

%=====================================================================
\section{Experimental Setup}\label{sec:setup}
%=====================================================================

\subsection{Network}
The network covers Tunis centre-ville, built with \texttt{netconvert} from an OpenStreetMap
extract dated 15 October 2026 (1{,}335 edges). Six junctions, J1--J6, were selected by a
deterministic rule over the OSM graph (degree $\geq3$, then signal tags, roundabout,
named, degree, major road, node id); they lie on Av.\ Habib Bourguiba, Av.\ Habib
Thameur, Rue de Rome and neighbouring streets. Signals were set at the six nodes, and other
junctions yield. Lane evidence is weak: only 3 of the 21 roads incident to the six
junctions carry an OSM lane-count tag, so the other 18 use SUMO defaults. Roads without a
speed tag use Tunisia's 50\,km/h urban limit. Bus geometry for 13 lines and 16 stops comes
from community-maintained data attributed to the operator TRANSTU, with assumed headways.

\subsection{Demand and scenarios}
No open traffic counts exist for these junctions, so demand is synthetic: trips generated
with SUMO's \texttt{randomTrips} from a calendar model at a base of 2{,}400 vehicles/h plus
pedestrians and buses. Twelve scenarios (peaks, Ramadan, rain, incidents, emergencies,
sensor faults) are defined; scenario events are hidden from the agents. This paper reports
three of them: \emph{morning peak}, \emph{emergency} (ambulances crossing the network) and
\emph{incident} (lane blockages). Each scenario was run with seeds 1--3 at a 1\,s step in
SUMO 1.27.

\subsection{Controllers and metrics}
We compare four controllers on identical demand files per seed: SUMO's \emph{fixed}-time
plans, SUMO's \emph{actuated} plans, \emph{traffix rules} (the full hierarchy with rule
fallbacks only) and \emph{traffix models} (the full hierarchy with trained models and
Max-Pressure lights). Metrics are mean vehicle waiting time, mean pedestrian waiting time,
emergency vehicle trip time, and the number of vehicle greens started while pedestrians
were still on a crossing. Differences are tested with paired $t$-tests over seeds.

%=====================================================================
\section{Results}\label{sec:results}
%=====================================================================

\begin{table*}[t]
\caption{Controller results on the Tunis network (mean of seeds 1--3, lower is better, best
in bold). Percentages compare traffix models with SUMO actuated. No difference is
significant at $p<0.05$.}
\label{tab:results}
\centering
\begin{tabular}{@{}llrrrr@{}}
\toprule
Scenario & Measure & SUMO fixed & SUMO actuated & traffix rules & traffix models \\
\midrule
Morning peak & Mean vehicle wait (s) & 122.6 & 118.5 & \textbf{101.8} & 105.9 ($-11\%$) \\
Emergency    & Mean vehicle wait (s) & \textbf{62.2} & 65.9 & 68.7 & 78.8 ($+20\%$) \\
Incident     & Mean vehicle wait (s) & \textbf{69.1} & 96.1 & 98.1 & 73.5 ($-24\%$) \\
\midrule
Morning peak & Mean pedestrian wait (s) & 140.2 & 144.1 & 128.1 & \textbf{116.3} ($-19\%$) \\
Emergency    & Mean pedestrian wait (s) & 139.1 & 152.2 & 107.6 & \textbf{98.3} ($-35\%$) \\
Incident     & Mean pedestrian wait (s) & 123.9 & \textbf{111.1} & 118.9 & 123.3 ($+11\%$) \\
\midrule
Emergency    & Emergency vehicle trip (s) & \textbf{197.0} & 213.2 & 232.3 & 240.0 ($+13\%$) \\
\midrule
All three    & Greens started during a crossing (per run) & 239--281 & 331--405 & 16--19 & \textbf{17--19} \\
\bottomrule
\end{tabular}
\end{table*}

Table~\ref{tab:results} gives the results.

\paragraph{Pedestrian safety} The clearest effect is on greens started while pedestrians
were still crossing: 331--405 per run with actuated plans and 239--281 with fixed plans,
against 16--19 with either traffix variant, a reduction of about 95\%. Because the rule-only
and model variants are equal here, the gain comes from the executor's clearance rule and not
from learning. It is not zero, so the pedestrian-safety objective is only partly met.

\paragraph{Waiting times} Pedestrian waiting fell against actuated in the morning peak
($-19\%$) and emergency ($-35\%$) scenarios and rose in the incident scenario ($+11\%$).
Vehicle waiting was mixed: lower in the morning peak ($-11\%$) and incident ($-24\%$)
scenarios, higher in the emergency scenario ($+20\%$). In two of three scenarios the
simplest baseline, fixed time, gave the lowest vehicle wait. With three seeds none of the
paired $t$-tests between traffix and actuated, or between models and rules, was
significant at $p<0.05$; the closest was pedestrian waiting in the morning peak, models
against rules, $-9\%$ ($p=0.06$).

\paragraph{Emergency vehicles} Emergency trips were slower with traffix (240.0\,s) than with
actuated (213.2\,s) or fixed plans (197.0\,s), even though the emergency agent requests
green ahead of the vehicle. We have not yet identified the cause; candidates are the
preparation horizon, the exit-full refusal rule and the extra vehicle queues in this
scenario. The emergency objective is not met.

\paragraph{Learned components} Seven of eleven models beat their rule on held-out data, for
example ambulance arrival time (mean error 43.7\,s against 92.6\,s for distance over speed)
and demand forecasting during Ramadan (0.026 against 0.103). The arrival model only matched
its rule, the pedestrian classifier is weak (F1 0.44), and the trained sensor-health
classifier lost to its rule (F1 0.75 against 0.98) and was removed. Random search over the
Max-Pressure parameters (8 trials) found nothing better than the defaults. After three
rounds of Fitted Q-Iteration, the RL light waited 231.4\,s (vehicle plus pedestrian) in the
final test against 221.9\,s for Max-Pressure, about 4\% worse, so Max-Pressure remains the
default.

\paragraph{Language layer} Every logged decision carries a reason code by construction, and in
our tests a false claim submitted to the audit endpoint is answered with an abstention. We
have not yet measured the explainer quantitatively (Section~\ref{sec:discussion}).

%=====================================================================
\section{Discussion and Threats to Validity}\label{sec:discussion}
%=====================================================================

We now examine our own protocol against the guidance of~\cite{riehl2026sumoitscontrol}, and
our design against~\cite{azevedo2016jade} and~\cite{li2026decoupled}.

\subsection{Evaluation protocol}
\emph{Single-seed tuning.} The Max-Pressure parameters were tuned by scoring each candidate
on one seed. Riehl et al.\ show that neighbouring parameter values can differ by less than
their run-to-run standard deviation (12.48 against 13.04 with standard deviations of about
2.3 over 20 runs, $p=0.46$), so a single run can rank them either way. Our conclusion that
``no candidate beat the defaults'' may therefore reflect noise. Candidates should be scored
by their mean over a fixed, documented set of seeds.

\emph{Too few replications.} Three seeds leave a paired $t$-test with little power, so our
``not significant'' results are weak evidence of no effect. The recommended range is 10--20
replications per configuration.

\emph{Multiple comparisons.} Our analysis script tests five controller pairs on three
metrics in each scenario, up to 180 tests over all twelve scenarios. At $\alpha=0.05$, about
nine would appear significant by chance even if no controller differed. A family-wise
correction such as Holm's~\cite{holm1979} is needed, and where normality is doubtful the
Wilcoxon signed-rank test~\cite{wilcoxon1945} should replace the $t$-test.

\emph{Simulation settings.} We measure from the first second, without a warm-up period, and
use a 1\,s step where 0.25\,s is recommended for urban scenes. Both may bias waiting times.

\emph{Baselines.} SUMO's default fixed and actuated plans are weak references, and fixed
time won two of our vehicle-wait comparisons. Following~\cite{azevedo2016jade}, a fixed plan
with time-of-day timings is a stronger and fairer baseline for learned controllers;
following~\cite{riehl2026sumoitscontrol}, an external Max-Pressure and a SCOOT/SCATS-style
group controller would show whether our gains, where present, come from the architecture or
from a weak baseline.

\emph{No maximum green.} Our Max-Pressure mode has no maximum green time, whereas the
flexible controller of~\cite{riehl2026sumoitscontrol} caps green at 50\,s and our own RL
agent at 60\,s. A busy approach can therefore keep green while a quiet one starves.

\emph{Fairness metrics.} Following~\cite{azevedo2016jade}, the standard deviation of trip and
waiting times should be reported next to the means, since controllers can trade a lower
average for more variable service.

\subsection{Architecture}
Azevedo et al.~\cite{azevedo2016jade} present centralised and distributed coordination as
alternatives. Our hierarchy combines them: decisions remain local and single-writer, while
the supermaster provides the network view through directives that expire. The
\emph{junction-only} controller in our experiment suite (masters without a supermaster)
measures what the coordinator adds; it was not part of the three-scenario run reported here.

\subsection{Language layer}
Li and Ke~\cite{li2026decoupled} found that a single repair round recovered many failed
LLM runs while further rounds did not. Our explainer abstains at the first failed citation
check; one retry with the error fed back is a cheap improvement. They also note that
smaller models break output schemas unless the output is validated, which argues for a
structured \texttt{\{answer, citations, abstained\}} output over citation extraction from
free text. Finally, their evaluation (task success, tokens and time to first result across
three difficulty levels) is a template for the benchmark Synapse lacks: a fixed set of
questions over saved logs, scored on citation correctness, correct abstention, cost and
latency, for both the LLM and the template engines.

\subsection{External validity}
Demand is synthetic, 18 of 21 junction roads have unverified lane counts, bus headways are
assumed and calendar effects are our own assumptions. The results describe this simulated
network only.

%=====================================================================
\section{Conclusion}
%=====================================================================

We built a six-junction multi-agent signal controller on a real street network in which
one master writes each light, control degrades predictably when data fails, and a language
layer explains decisions from a log it cannot change. The architecture rules are enforced
mechanically, and every decision is logged with a reason. The clearest measured benefit is
a 95\% reduction in greens started during pedestrian crossings, which comes from the safety
rules rather than from learning. Effects on waiting times are mixed and not yet
significant, and emergency vehicles were slower than with SUMO's own plans. Our next steps
follow from the critique above: tune over several seeds, run 10--20 seeds over all twelve
scenarios with corrected tests and a warm-up period, add a maximum green, compare against
time-of-day and external adaptive baselines, find the cause of the emergency slowdown, and
benchmark the language layer.

\section*{Data and Code Availability}
Code, network, scenario definitions and the scripts that produce every number in this
paper are in the project repository: \texttt{scripts/run\_experiments.py} and
\texttt{scripts/analyze.py} for Table~\ref{tab:results}, \texttt{models/metrics.json} for the
model comparisons, and \texttt{scripts/check\_architecture.py} for the architecture rules.

\bibliographystyle{IEEEtran}
\bibliography{references}

\end{document}
